\section{Experiment 1: POPE Object Hallucination Detection}

\subsection{Results}

\Cref{tab:pope_results} reports all evaluation metrics
for the four models across the three POPE splits.
Columns are grouped by split,
each containing Accuracy, Precision, Recall, F1, Yes Ratio,
Hallucination Rate (HR),
and the four confusion matrix entries (TP/TN/FP/FN).

\begin{table}[htbp]
  \centering
  \caption{Comprehensive POPE results across three splits.
    Best values per metric are highlighted in pink.}
  \label{tab:pope_results}
  \small
  \setlength{\tabcolsep}{4pt}
  \begin{tabular}{ll cccc}
    \toprule
    Split & Metric
      & GPT-5.4-mini & Gemini 2.5 Flash & Qwen3.5-35B & Qwen3-VL-235B \\
    \midrule
    \multirow{7}{*}{Random}
      & Accuracy  & 0.9010 & 0.9160 & \cellcolor{pink!25}0.9260 & 0.9200 \\
      & Precision & 0.9831 & 0.9814 & 0.9754 & \cellcolor{pink!25}0.9929 \\
      & Recall    & 0.8156 & 0.8477 & \cellcolor{pink!25}0.8737 & 0.8457 \\
      & F1        & 0.8916 & 0.9097 & \cellcolor{pink!25}0.9218 & 0.9134 \\
      & Yes Ratio & 0.414  & 0.431  & 0.447  & 0.425 \\
      & HR        & 0.70\% & 0.80\% & 1.10\% & \cellcolor{pink!25}0.30\% \\
      & TP/TN/FP/FN & 407/494/7/92 & 423/493/8/76
        & 436/490/11/63 & 422/498/3/77 \\
    \midrule
    \multirow{7}{*}{Popular}
      & Accuracy  & 0.8570 & 0.8620 & 0.8730 & \cellcolor{pink!25}0.8970 \\
      & Precision & 0.8836 & 0.8800 & 0.8750 & \cellcolor{pink!25}0.9459 \\
      & Recall    & 0.8216 & 0.8377 & \cellcolor{pink!25}0.8697 & 0.8417 \\
      & F1        & 0.8515 & 0.8583 & 0.8724 & \cellcolor{pink!25}0.8908 \\
      & Yes Ratio & 0.464  & 0.475  & 0.496  & 0.444 \\
      & HR        & 5.40\% & 5.70\% & 6.20\% & \cellcolor{pink!25}2.40\% \\
      & TP/TN/FP/FN & 410/447/54/89 & 418/444/57/81
        & 434/439/62/65 & 420/477/24/79 \\
    \midrule
    \multirow{7}{*}{Adversarial}
      & Accuracy  & 0.8280 & 0.8730 & 0.8520 & \cellcolor{pink!25}0.8750 \\
      & Precision & 0.8385 & 0.8924 & 0.8318 & \cellcolor{pink!25}0.8979 \\
      & Recall    & 0.8116 & 0.8477 & \cellcolor{pink!25}0.8818 & 0.8457 \\
      & F1        & 0.8248 & 0.8695 & 0.8560 & \cellcolor{pink!25}0.8710 \\
      & Yes Ratio & 0.483  & 0.474  & 0.529  & 0.470 \\
      & HR        & 7.80\% & 5.10\% & 8.90\% & \cellcolor{pink!25}4.80\% \\
      & TP/TN/FP/FN & 405/423/78/94 & 423/450/51/76
        & 440/412/89/59 & 422/453/48/77 \\
    \bottomrule
  \end{tabular}
\end{table}

\begin{figure}[htbp]
  \centering
  \includegraphics[width=0.75\textwidth]{fig1_pope_hr_bar.pdf}
  \caption{Hallucination rates of four MLLMs
    across the three POPE splits
    (Random, Popular, Adversarial).}
  \label{fig:pope_hr_bar}
\end{figure}

\subsection{Analysis}

\subsubsection{Three Splits Probe Three Error Mechanisms}
\label{sec:pope_difficulty}

\textbf{Monotonic difficulty gradient.}
From Random to Popular to Adversarial,
the Hallucination Rate of all four models increases monotonically
(\Cref{tab:pope_results,fig:pope_hr_bar}),
confirming the validity of POPE's three-tier negative-sampling design.
Beyond a difficulty ranking,
however, the three splits also expose
three qualitatively distinct error mechanisms
that map cleanly onto the faithfulness taxonomy.

\textbf{Random split: object-level perception.}
Negative objects are sampled uniformly,
so most candidates are visually unrelated to the scene
(e.g., asking about a car in a kitchen photo).
All four models achieve HR between 0.30\% and 1.10\%,
indicating that the baseline ability to reject obviously absent objects
is near saturation.
This split therefore measures raw object-level visual perception,
with negligible interference from language priors.

\textbf{Popular split: resistance to linguistic priors.}
Negative objects are high-frequency co-occurring categories
that are statistically associated with the scene
but absent from the specific image
(e.g., asking about a backpack in an outdoor eating scene).
HR rises sharply to 2.40\%--6.20\%,
roughly 5--8 times the Random level.
The surge isolates the contribution of language-prior bias:
when the visual evidence is ambiguous,
models fall back on the statistical co-occurrence patterns
learned during pretraining.
Case~2 in \Cref{sec:case2_pope} provides a direct instance of this mechanism.

\textbf{Adversarial split: fine-grained semantic discrimination.}
Negative objects are semantically similar to objects that do appear in the image
(e.g., asking about a bear when only a teddy bear is present).
HR rises further to 4.80\%--8.90\%
and the model can no longer rely on scene-level associations:
the visual encoder must resolve fine-grained semantic boundaries.
Case~1 in \Cref{sec:case1_pope}
shows that this is the hardest of the three failure modes,
because the underlying cause, confusion of conceptual hierarchies, is
not addressable by scale alone.

In sum,
the three splits jointly diagnose
where in the perception--prior pipeline
an object-level faithfulness hallucination originates,
rather than merely how often it occurs.

\subsubsection{Cross-Model Comparison}
\label{sec:pope_cross_model_comparison}

\textbf{Within the Qwen family, visual-encoder capacity dominates.}
Qwen3-VL-235B achieves the lowest HR
on every split (0.30\%, 2.40\%, 4.80\%)
and the highest F1 on Popular and Adversarial.
Compared with Qwen3.5-35B-A3B,
which shares the same training family
but uses a smaller activated visual encoder
(3\,B vs.\ 22\,B activated parameters),
HR is consistently 50--75\% lower
despite identical instruction-tuning regimes.
Holding the language backbone family fixed,
visual-encoder scale is therefore the dominant factor
within this family;
whether the same holds across families cannot be established here,
because the closed-source dense models
do not disclose visual-encoder parameters.

\textbf{GPT-5.4-mini: conservative strategy with asymmetric costs.}
GPT-5.4-mini has the lowest Yes Ratio of all models (0.414 on Random)
and the highest Precision (0.9831 on Random).
This conservative profile yields a low FP count
but the highest FN count (92 on Random),
so its hallucinations manifest predominantly
as omissions rather than fabrications.
In safety-sensitive applications
this asymmetric cost may be desirable
(it is generally preferable not to invent objects),
but it caps the achievable F1:
HR on Adversarial rises more than tenfold relative to Random
(0.70\% $\to$ 7.80\%)
once the Yes/No decision can no longer be safely defaulted to ``No.''

\textbf{Gemini 2.5 Flash: atypical Popular--Adversarial inversion.}
Gemini is the only model whose HR \emph{decreases}
from Popular (5.70\%) to Adversarial (5.10\%),
though the gap is small (0.6\,pp).
Since the standard interpretation
(Adversarial is harder than Popular)
holds for the other three models,
this inversion suggests
that Gemini's visual perception is comparatively robust
to fine-grained semantic distractors
but relatively sensitive to scene-level
co-occurrence priors.
A plausible mechanism
is a perception pipeline
that grounds responses more tightly in visible features
than in scene context,
which would also be consistent with
its moderate but stable Popular performance.

\textbf{Qwen3.5-35B: largest split-difficulty slope.}
Qwen3.5-35B-A3B exhibits the steepest HR increase
across splits (1.10\% $\to$ 6.20\% $\to$ 8.90\%, a 7.8\,pp rise),
the largest in the panel.
With only 3\,B activated visual parameters,
the visual encoder is insufficient
to support fine-grained discrimination
once the negative samples become semantically adjacent.
The high Recall (highest in three of three splits)
reflects an opposite pressure
toward answering ``Yes,''
which is also visible in the panel's highest Yes Ratio (0.529 on Adversarial).

\textbf{Synthesis.}
Whereas the previous subsection diagnosed
where in the perception--prior pipeline an error originates,
the cross-model view identifies
who is vulnerable to each mechanism.
The comparison highlights two axes:
visual-encoder capacity,
which explains the gap between the two Qwen models
and the rough position of Gemini;
and response-bias calibration,
which separates the conservative profile of GPT-5.4-mini
from the more permissive profiles of the Qwen models.
Models near the corner of high capacity and balanced calibration
(Qwen3-VL) dominate the leaderboard;
modest capacity with permissive calibration
(Qwen3.5) sacrifices precision under adversarial pressure;
and modest capacity with conservative calibration
(GPT-5.4-mini) sacrifices recall.
This two-axis view will recur in
\Cref{sec:mathvista,sec:vqarad}
under different task pressures.

\subsection{Failure Case Analysis}

\subsubsection{Case 1: Semantic Granularity Confusion (Adversarial)}
\label{sec:case1_pope}

This case instantiates
the fine-grained semantic discrimination mechanism
identified in \Cref{sec:pope_difficulty}.
\Cref{tab:pope_case1,fig:pope_case1} present a sample from the Adversarial split.
The image contains
a teddy bear, a person, and a book.
The question asks whether there is a bear in the image,
and the ground truth is No.
All four models answer Yes,
indicating a systematic error rather than an individual model deficiency.

\begin{table}[htbp]
  \centering
  \begin{minipage}[c]{0.5\textwidth}
    \centering
    \caption{POPE Case 1: sample information and model responses.}
    \label{tab:pope_case1}
    \small
    \begin{tabular}{@{}ll@{}}
      \toprule
      \multicolumn{2}{c}{\textbf{Sample Information}} \\
      \midrule
      Image & COCO\_val2014\_000000567886 \\
      Objects present & teddy bear, person, book \\
      Objects absent & bear, car, tv \\
      Question & Is there a bear in the image? \\
      Ground Truth & No \\
      \midrule
      \multicolumn{2}{c}{\textbf{Model Responses}} \\
      \midrule
      GPT-5.4-mini & Yes \\
      Gemini 2.5 Flash & Yes \\
      Qwen3.5-35B-A3B & Yes \\
      Qwen3-VL-235B & Yes \\
      \bottomrule
    \end{tabular}
  \end{minipage}%
  \hspace{0.02\textwidth}%
  \begin{minipage}[c]{0.45\textwidth}
    \centering
    \includegraphics[width=0.9\textwidth]{pope_000000567886.jpg}
    \captionof{figure}{POPE Case 1: a teddy bear on a couch.}
    \label{fig:pope_case1}
  \end{minipage}
\end{table}

The root cause is semantic granularity confusion.
The image contains a teddy bear,
an independent category in the COCO label set
\citep{linMicrosoftCOCOCommon2015},
while the question asks about a bear (a real animal).
All four models associate ``teddy bear'' with ``bear''
at the semantic level,
triggering a positive judgment upon seeing the toy.
This is essentially conceptual hierarchy confusion:
models conflate visually similar objects with semantically identical ones.
This error cannot be resolved
by increasing model scale or improving visual encoder resolution,
as it concerns the clarity of conceptual boundaries
in semantic representations
rather than the discriminability of low-level visual features.

\subsubsection{Case 2: Linguistic Prior Bias (Popular)}
\label{sec:case2_pope}

This case exemplifies
the linguistic-prior bias mechanism
isolated by the Popular split (\Cref{sec:pope_difficulty}).
\Cref{tab:pope_case2,fig:pope_case2} present a sample from the Popular split.
The image shows
a person eating a sandwich and a hot dog outdoors.
The question asks whether there is a backpack in the image,
and the ground truth is No.
Again, all four models answer Yes.

\begin{table}[htbp]
  \centering
  \begin{minipage}[c]{0.6\textwidth}
    \centering
    \caption{POPE Case 2: sample information and model responses.}
    \label{tab:pope_case2}
    \small
    \begin{tabular}{@{}ll@{}}
      \toprule
      \multicolumn{2}{c}{\textbf{Sample Information}} \\
      \midrule
      Image & COCO\_val2014\_000000569839 \\
      Objects present & person, sandwich, hot dog \\
      Objects absent & car, bottle, backpack \\
      Question & Is there a backpack in the image? \\
      Ground Truth & No \\
      \midrule
      \multicolumn{2}{c}{\textbf{Model Responses}} \\
      \midrule
      GPT-5.4-mini & Yes \\
      Gemini 2.5 Flash & Yes \\
      Qwen3.5-35B-A3B & Yes \\
      Qwen3-VL-235B & Yes \\
      \bottomrule
    \end{tabular}
  \end{minipage}%
  \hspace{0.02\textwidth}%
  \begin{minipage}[c]{0.3\textwidth}
    \centering
    \includegraphics[width=0.9\textwidth]{pope_000000569839.jpg}
    \captionof{figure}{POPE Case 2: an outdoor dining scene.}
    \label{fig:pope_case2}
  \end{minipage}
\end{table}

The root cause is linguistic prior bias.
In outdoor scenes with people eating,
``backpack'' is a high-frequency co-occurring object.
The statistical prior in the language module
(outdoor + person + eating $\to$ high backpack probability)
overrides the visual module's actual perception,
causing the model to produce an affirmative answer.
This is precisely the mechanism that the Popular split exploits:
models fill visual uncertainty
with linguistic priors from training data.
The contrast between the two cases further confirms
that the Popular and Adversarial splits
target distinct error mechanisms,
with Popular primarily exploiting linguistic priors
and Adversarial primarily testing semantic discrimination.
